% Created 2024-12-10 mar 16:09
% Intended LaTeX compiler: pdflatex
\documentclass[presentation]{beamer}
\usepackage[utf8]{inputenc}
\usepackage[T1]{fontenc}
\usepackage{graphicx}
\usepackage{longtable}
\usepackage{wrapfig}
\usepackage{rotating}
\usepackage[normalem]{ulem}
\usepackage{amsmath}
\usepackage{amssymb}
\usepackage{capt-of}
\usepackage{hyperref}
\usepackage{minted}
\usetheme{metropolis}
\usetheme{default}
\author{Tu Nombre}
\date{\today{}}
\title{Bases de Datos en Big Data: SQL vs NoSQL}
\hypersetup{
 pdfauthor={Tu Nombre},
 pdftitle={Bases de Datos en Big Data: SQL vs NoSQL},
 pdfkeywords={},
 pdfsubject={},
 pdfcreator={Emacs 27.1 (Org mode 9.6.18)}, 
 pdflang={Spanish}}
\begin{document}

\maketitle

\begin{frame}[label={sec:org0f02cd2}]{Introducción}
Las bases de datos son el núcleo de los sistemas Big Data:
\begin{itemize}
\item Almacenan, organizan y recuperan grandes volúmenes de información.
\item Evolución hacia dos paradigmas principales:
\begin{itemize}
\item \alert{\alert{SQL (relacionales)}}.
\item \alert{\alert{NoSQL (no relacionales)}}.
\end{itemize}
\end{itemize}

---
\end{frame}

\begin{frame}[label={sec:org97f129f}]{Bases de Datos Relacionales (SQL)}
\begin{block}{Características Clave}
\begin{enumerate}
\item \alert{\alert{Estructura Definida}}:
\begin{itemize}
\item Esquema predefinido con tablas y relaciones.
\end{itemize}
\item \alert{\alert{Consistencia ACID}}:
\begin{itemize}
\item Atomicidad, Consistencia, Aislamiento, Durabilidad.
\end{itemize}
\item \alert{\alert{Consulta Eficiente}}:
\begin{itemize}
\item SQL permite consultas complejas sobre grandes volúmenes.
\end{itemize}
\end{enumerate}
\end{block}

\begin{block}{Ventajas}
\begin{itemize}
\item Alta integridad de datos.
\item Estandarización en la industria.
\item Excelente para relaciones complejas (ERP, CRM).
\end{itemize}
\end{block}

\begin{block}{Desventajas}
\begin{itemize}
\item Rigidez ante datos no estructurados.
\item Escalabilidad limitada (vertical).
\item No optimizadas para Big Data.
\end{itemize}
\end{block}

\begin{block}{Ejemplos}
\begin{itemize}
\item \alert{\alert{MySQL}}: Open source.
\item \alert{\alert{PostgreSQL}}: Avanzada y flexible.
\item \alert{\alert{Oracle Database}}: Enfocada en aplicaciones empresariales.
\item \alert{\alert{Microsoft SQL Server}}: Con herramientas de análisis.
\end{itemize}

---
\end{block}
\end{frame}

\begin{frame}[label={sec:orgfdf0edb}]{Bases de Datos NoSQL}
\begin{block}{Características Clave}
\begin{enumerate}
\item \alert{\alert{Flexibilidad en el Modelo de Datos}}:
\begin{itemize}
\item Datos sin esquema predefinido (JSON, documentos, columnas).
\end{itemize}
\item \alert{\alert{Escalabilidad Horizontal}}:
\begin{itemize}
\item Distribución en múltiples nodos.
\end{itemize}
\item \alert{\alert{Tipos de Modelos}}:
\begin{itemize}
\item Clave-valor (Redis), Documentales (MongoDB), Columnas (Cassandra), Grafos (Neo4j).
\end{itemize}
\end{enumerate}
\end{block}

\begin{block}{Ventajas}
\begin{itemize}
\item Adaptabilidad a datos no estructurados.
\item Escalabilidad económica y efectiva.
\item Alto rendimiento para aplicaciones en tiempo real.
\end{itemize}
\end{block}

\begin{block}{Desventajas}
\begin{itemize}
\item Consistencia relajada (modelo CAP).
\item Menos ideales para relaciones complejas.
\item Curva de aprendizaje más pronunciada.
\end{itemize}
\end{block}

\begin{block}{Ejemplos}
\begin{itemize}
\item \alert{\alert{MongoDB}}: Basada en documentos.
\item \alert{\alert{Cassandra}}: Para grandes volúmenes distribuidos.
\item \alert{\alert{Redis}}: En memoria, para alta velocidad.
\item \alert{\alert{Neo4j}}: Optimizada para grafos.
\end{itemize}

---
\end{block}
\end{frame}

\begin{frame}[label={sec:org7be2788}]{Comparación: SQL vs NoSQL}
$\backslash$[
\begin{array}{|l|l|l|}
\hline
\textbf{Aspecto}      & \textbf{SQL}                     & \textbf{NoSQL}                         \\
\hline
Estructura            & Esquema fijo, tablas.           & Flexible, sin esquema.                 \\
\hline
Consistencia          & ACID (total).                  & Relajada (eventual en algunos casos).  \\
\hline
Escalabilidad         & Vertical.                      & Horizontal.                            \\
\hline
Tipo de Datos         & Estructurados.                 & Semi/no estructurados.                 \\
\hline
Relaciones            & Ideal para relaciones.         & Menos optimizado para relaciones.      \\
\hline
Ejemplos              & MySQL, PostgreSQL.             & MongoDB, Cassandra.                    \\
\hline
\end{array}
$\backslash$]

---
\end{frame}

\begin{frame}[label={sec:org19ca14a}]{Bases de Datos en Big Data}
\begin{block}{Principales Características}
\begin{enumerate}
\item \alert{\alert{Alta Escalabilidad}}:
\begin{itemize}
\item Crecimiento horizontal para manejar volúmenes masivos.
\end{itemize}
\item \alert{\alert{Tolerancia a Fallos}}:
\begin{itemize}
\item Replicación de datos entre nodos.
\end{itemize}
\item \alert{\alert{Datos No Estructurados}}:
\begin{itemize}
\item Logs, multimedia, JSON, etc.
\end{itemize}
\item \alert{\alert{Procesamiento Paralelo}}:
\begin{itemize}
\item Integración con Hadoop y Spark.
\end{itemize}
\end{enumerate}

---
\end{block}
\end{frame}

\begin{frame}[label={sec:org1a4e963}]{¿Cómo Elegir?}
\begin{block}{Usar SQL si:}
\begin{itemize}
\item Necesitas transacciones consistentes.
\item Trabajas con datos estructurados y relaciones complejas.
\item Prefieres herramientas estandarizadas.
\end{itemize}
\end{block}

\begin{block}{Usar NoSQL si:}
\begin{itemize}
\item Manejas datos no estructurados o semi-estructurados.
\item Requieres escalabilidad horizontal.
\item Trabajas con redes sociales, IoT o análisis en tiempo real.
\end{itemize}

---
\end{block}
\end{frame}

\begin{frame}[label={sec:org74b86f7}]{Conclusión}
SQL y NoSQL son complementarios:
\begin{itemize}
\item \alert{\alert{SQL}}: Para datos estructurados y relaciones complejas.
\item \alert{\alert{NoSQL}}: Para manejar volumen, variedad y velocidad en Big Data.
\item La elección depende de los requisitos de la aplicación y el tipo de datos.
\end{itemize}
\end{frame}
\end{document}
